\documentclass[11pt]{article}
\usepackage{fontspec}
\setmainfont{Latin Modern Roman}
\usepackage{amsmath,amssymb,mathtools,unicode-math}
\setmathfont{Latin Modern Math}
\usepackage{xurl,hyperref}
\hypersetup{colorlinks=true,urlcolor=blue,pdftitle={Problem statement}}
\newcommand{\sref}[2]{\hyperlink{source:#1:#2}{[S#2]}}
\newcommand{\cataloguescope}[1]{\paragraph{Mathematical question.}#1}
\begin{document}
\section{Algebraic Stories — Generic $k$-rank}\label{statement}

\subsection{Definitions and mathematical statement}

\paragraph{Shared notation.}$\mathbb N=\{1,2,\ldots\}$, and $\mathbb Z,\mathbb Q,\mathbb R,\mathbb C$ denote the integers, rationals, reals and complex numbers. Any other conventions are specified locally.


Work over $\mathbb C$. For $n\geq1$, let $S=\mathbb C[x_1,\ldots,x_n]$, with $S_e$ its vector space of homogeneous polynomials of total degree $e$. For $k,d\geq1$ and $f\in S_{kd}$, define
\[\operatorname{rk}_k(f)=\min\left\{s\geq0:f=\sum_{j=1}^{s}g_j^k,\quad g_j\in S_d\right\}.
\]
A general form means one whose coefficient vector lies in a suitable nonempty Zariski open subset of $S_{kd}$. Write $\operatorname{rk}_k^\circ(kd,n)$ for the value on general forms.
\paragraph{Statement recorded in the supplied source.}
For every triple of positive integers $(k,d,n)$, determine $\operatorname{rk}_k^\circ(kd,n)$. \sref{statement}{2}
\subsection{Short English statement}
Determine how many powers of degree-$d$ complex forms are needed for a general degree-$kd$ form in any number of variables.
\subsection{Sources}
\begin{description}
\item[\hypertarget{source:statement:1}{S1}] ULAM.AI and the UnsolvedMath Contributors, UnsolvedMath: A Curated Collection of Open Mathematics Problems, catalogue entry. CC BY 4.0. \url{https://huggingface.co/datasets/ulamai/UnsolvedMath}.
\item[\hypertarget{source:statement:2}{S2}] R. Fröberg, S. Lundqvist, A. Oneto and B. Shapiro, \emph{Algebraic Stories from One and from the Other Pockets}, arXiv:1801.01692 (2018), §1.1, Problem A. \url{https://arxiv.org/abs/1801.01692}.
\end{description}
\label{end:statement}
\end{document}
